\documentclass{article}
\usepackage{enumitem}
\usepackage{amsmath}
\begin{document}

\title{Chapter 3: Plant Kingdom}
\date{}
\maketitle

\begin{enumerate}[label=Q\arabic*.]

\item In the plant kingdom, which group is considered the most primitive land plants?
\\(A) Pteridophytes
\\(B) Bryophytes
\\(C) Gymnosperms
\\(D) Angiosperms

\item The dominant phase in the life cycle of bryophytes is:
\\(A) Sporophyte
\\(B) Gametophyte
\\(C) Both equally dominant
\\(D) Sporophyte in mosses, gametophyte in liverworts

\item Which of the following is a correct statement about pteridophytes?
\\(A) They produce seeds enclosed in fruits
\\(B) They are the first vascular plants with independent sporophyte
\\(C) They have dominant gametophyte generation
\\(D) They reproduce by pollen grains

\item In gymnosperms, the female gametophyte is represented by:
\\(A) Ovule
\\(B) Megaspore mother cell
\\(C) Endosperm
\\(D) Both ovule and megaspore

\item Which of the following features is unique to angiosperms among seed plants?
\\(A) Presence of ovules
\\(B) Double fertilization leading to triploid endosperm
\\(C) Presence of pollen grains
\\(D) Alternation of generations

\item Assertion: Bryophytes are called amphibians of the plant kingdom.\\
Reason: They require water for fertilization but can survive on land.
\\(A) Both true, Reason is correct explanation
\\(B) Both true, Reason is not correct explanation
\\(C) Assertion true, Reason false
\\(D) Assertion false, Reason true

\item The prothallus in pteridophytes represents which generation?
\\(A) Sporophyte
\\(B) Gametophyte
\\(C) Both sporophyte and gametophyte
\\(D) Neither; it is a transitional form

\item Heterospory in pteridophytes is significant because it is:
\\(A) A precursor to seed habit in higher plants
\\(B) A mechanism for vegetative reproduction
\\(C) Essential for nitrogen fixation
\\(D) Required for formation of spore coat

\item In \textit{Pinus}, the microsporophyll bears:
\\(A) Ovules
\\(B) Microsporangia containing pollen grains
\\(C) Archegonia
\\(D) Megasporophylls

\item Which group of plants shows the phenomenon of polyembryony as a natural occurrence?
\\(A) Angiosperms exclusively
\\(B) Gymnosperms, especially conifers
\\(C) Pteridophytes
\\(D) Bryophytes

\item The w-type of alternation of generations, where sporophyte is completely dependent on gametophyte, is seen in:
\\(A) Ferns
\\(B) Mosses
\\(C) Gymnosperms
\\(D) Angiosperms

\item Match the following plant groups with their key features:\\
1. Bryophytes $\rightarrow$ (a) No vascular tissue, dominant gametophyte\\
2. Pteridophytes $\rightarrow$ (b) Vascular tissue, no seeds\\
3. Gymnosperms $\rightarrow$ (c) Naked seeds, dominant sporophyte\\
4. Angiosperms $\rightarrow$ (d) Seeds enclosed in fruit
\\(A) 1-a, 2-b, 3-c, 4-d
\\(B) 1-b, 2-a, 3-d, 4-c
\\(C) 1-c, 2-d, 3-a, 4-b
\\(D) 1-d, 2-c, 3-b, 4-a

\item Which of the following is NOT a characteristic of gymnosperms?
\\(A) Ovules borne on megasporophylls
\\(B) Presence of fruit enclosing seeds
\\(C) Pollen grains transported by wind
\\(D) Dominant sporophyte generation

\item Isogamy, anisogamy, and oogamy represent:
\\(A) Types of asexual reproduction
\\(B) Progressive differentiation in size and form of gametes
\\(C) Types of vegetative propagation
\\(D) Different stages of meiosis

\item Selaginella and Salvinia are placed in pteridophytes because they:
\\(A) Produce seeds after fertilization
\\(B) Are vascular plants producing spores without seeds
\\(C) Have dominant gametophyte
\\(D) Require water for photosynthesis only

\end{enumerate}
\end{document}
